% ==========================================================================
% GENERATED by scripts/split_paper.py. Do not edit this file.
%
% Source: imported_papers/probe-norm/content/body.tex, content/supplement.tex
% Part 8: The Trained Modules
%
% Edits belong upstream, in the published paper. To change what the
% chapter includes, change imports/bounded.toml and re-run
% the splitter.
% ==========================================================================
\section{The Trained Modules}
\label{app:trained}
\label{app:impl}
All quantities are computed per checkpoint before aggregation, and pooled medians include every checkpoint in the stated group.

\paragraph{Encoder.} All modules share the encoder architecture and differ only in their dynamics. Each module learns its own encoder, which receives the current and the preceding $128\times128$ frame. Three stride-2 convolutional blocks with group normalization and GELU are followed by a two-layer MLP whose output has $D=256$ coordinates and ends in layer normalization~\citep{Ba2016LayerNorm}. This final normalization sets each encoding's norm from its direction, which gives the analytic reference radius of Section~\ref{app:radius}.

\paragraph{Non-Linear.} The residual network $F_\theta(\z,\mathbf a)=\z+g([\z,\mathbf a])$, where $g$ is an MLP with two hidden layers of GELU units, each as wide as the latent. Nothing in this module restricts the norm of its output.

\paragraph{Linear.} The linear map $F_\theta(\z,\mathbf a)=\K\z+\B\mathbf a$, with $\K$ initialized near the identity and $\B$ at zero. Repeated application of $\K$ can grow or shrink the norm without limit.

\paragraph{Linear-Penalized.} The same linear map, trained with the additional penalty $\lambda_s\max(0,s(\K)-1)$, where $s(\K)$ is the spectral radius. The penalty pushes $s(\K)$ toward at most one but does not enforce it, and even $s(\K)\le1$ does not bound the amplification over a finite horizon. Linear-Penalized on cartpole has spectral radius $1.02$, but its operator powers reach norm $88$.

\paragraph{Orthogonal.} The linear map with $\K$ orthogonal by construction, through the Cayley transform $\K=(I+A)^{-1}(I-A)$ of the skew-symmetric matrix $A=W-W^\top$~\citep{Helfrich2018}. We initialize $W$ from a least-squares fit to consecutive encodings~\citep{Schmid2010}. $\K$ preserves the norm exactly, so under bounded actions the norm can grow at most linearly with the horizon. It still need not stay at the reference radius.
\label{app:dmd}

These constraints follow work on stable latent dynamics~\citep{Lusch2018,Morton2018,KolterManek2019}. There they serve to improve the dynamics~\citep{Asadi2018,Bevanda2021,Erichson2019}, and here they vary how the norm evolves over a rollout.

\paragraph{Training variants.} Each variant changes one component of the base models. The squared-error variant replaces the cosine prediction loss. The fixed-radius variant rescales every predicted step to the norm of the initial encoding. The SimNorm variant applies SimNorm with $32$ groups of eight coordinates to encoder and dynamics outputs~\citep{Hansen2023}. The last variant removes the encoder's final layer normalization.

\paragraph{Training objective.} The objective combines prediction and representation regularization:
\begin{equation}
    \mathcal{L} = \lambda_p\mathcal{L}_\text{pred}
    + \lambda_v\mathcal{L}_\text{var}(Z_{\rm enc})
    + \lambda_v\lambda_c\mathcal{L}_\text{cov}(Z_{\rm enc}),
    \label{eq:base_loss}
\end{equation}
where $Z_{\rm enc}$ contains the online encoder outputs for the observed training frames. The term $\mathcal L_{\mathrm{pred}}$ averages cosine or squared error over $k=1,\dots,6$ against moving-average encoder targets. The variance and covariance penalties apply only to the encodings, and predicted latents enter only $\mathcal L_{\mathrm{pred}}$. Following VICReg~\citep{Bardes2022}, the variance term discourages constant representations, and the covariance term discourages redundant latent coordinates. We mask a quarter of image patches during training.

\paragraph{Hyperparameters.} Random searches of $32$ trials select the hyperparameters, first the encoder architecture on the Non-Linear module, then the optimization and loss weights for each module and environment. Both stages score a trial by the probe's held-out one-step state error. Acrobot reuses the cartpole values, and all variants keep the values of the base models without retuning (Table~\ref{tab:hyperparams}). Each run trains for $2000$ epochs, and we keep the epoch between $100$ and $2000$ with the lowest one-step validation probe error.

\paragraph{Data and probes.} Images are rendered at $128\times128$. State targets contain angle sines and cosines, velocities, and cart position where applicable. Episodes use uniformly random actions. All checkpoints share $400$ held-out episodes, $300$ for reference and $100$ for scoring. Probes train on $240$ reference episodes for $300$ epochs with AdamW and standardized targets, and the remaining $60$ give held-out $R^2$. Errors use outputs transformed back to state units.

\input{tables/table_hyperparams.tex}

\subsection{Checks on learned encoders}
\label{app:encoder}
Before evaluating predictions, we check the probes on exact encodings, where no dynamics error is present. We measure how much the probe's output changes when an encoding is rescaled without changing its state, which sets how strongly norm drift can distort the score, and whether an encoding's norm carries state that its direction lacks, which re-projection would discard.

\paragraph{Measurements.} \emph{Scale sensitivity} is the increase in a probe's error when exact encodings are multiplied by $c$, relative to $c=1$, divided by the reference mean's error. The \emph{radius check} compares probes fitted on directions alone with probes fitted on directions and norms. We standardize the inputs per coordinate, because without standardization, probes on unnormalized encoders can become nearly constant ($R^2$ as low as $-0.005$, compared with at least $0.64$ after standardization).

\paragraph{Probe fits.} Nearest-encoding factors, radius checks, and held-out $R^2$ average three probe fits. Other factors, encoding-rescaling tests, and probe variants use the first fit. We exclude first-fit results with $R^2\le0.5$, which removes six of $360$ checkpoints, four of them without final encoder normalization.

\input{tables/table_scale.tex}
\paragraph{Results.} Rescaling exact encodings to the upper edge of the off-radius band already raises the probe's error, and doubling their norm adds more than the reference mean's error (Table~\ref{tab:scale}). Linear, ReLU, and GELU probes remain scale-sensitive at every tested capacity, whereas tanh probes become less sensitive once their units saturate (Table~\ref{tab:v2_capacity}). The radius check finds no state information in the norm beyond the direction (Table~\ref{tab:arms}), so re-projection can be applied to these encoders without losing state.

\paragraph{Reference radius.}
\label{app:radius}
For layer normalization $\z=\boldsymbol\gamma\odot\hat{\mathbf x}+\boldsymbol\beta$, we solve $\|(c\mathbf u-\boldsymbol\beta)\oslash\boldsymbol\gamma\|^2=D$ along unit direction $\mathbf u$. With nonzero $\gamma_i$, let $a=\sum_i u_i^2/\gamma_i^2$, $b=\sum_i u_i\beta_i/\gamma_i^2$, and $k=\sum_i\beta_i^2/\gamma_i^2$. For $k<D$, the unique positive solution is
\begin{equation}
c_{\mathrm{LN}}(\mathbf u)=\frac{b+\sqrt{b^2+a(D-k)}}{a}.
\end{equation}
All $240$ layer-normalized encoders satisfy these conditions. The formula neglects the stabilizer $\epsilon$. At input variance $v$, the normalized squared norm is $Dv/(v+\epsilon)$ instead of $D$. Rescaling an arbitrary prediction does not enforce the additional zero-mean constraint and need not produce an encoder output. For SimNorm, each coordinate group sums to one, so $c=1/\sum_{i\in g}u_i$ for any group $g$ along a valid direction. Re-projection leaves SimNorm outputs unchanged.

\input{tables/table_capacity.tex}

\subsection{Controlled comparisons}
\label{app:full_results}
\label{app:trained_more}
Table~\ref{tab:arms_full} reports medians over ten seeds per module for the base models and five in each variant. We examine alternative radius estimates, probe fitting on rescaled encodings, and removal of encoder normalization.

\input{tables/table_arms_full.tex}

\input{tables/table_readouts_full.tex}

\paragraph{Norm growth and error.} Tables~\ref{tab:arms_full} and~\ref{tab:readouts_full} report the results by training setting, architecture, and task. Confidence intervals use $2000$ bootstrap resamples within environment--module groups. Adding norm to direction in probes fitted on predictions improves held-out $R^2$ by at most $0.018$ at $T=5$ and $0.066$ at $T=20$. At $T=20$ the gain exceeds the protocol's threshold of $0.05$ on one of $180$ checkpoints, on cartpole, where no module beats persistence at that horizon.

The analytic radius $c_{\mathrm{LN}}(\mathbf u)$ differs from the nearest-encoding estimate by a median $0.2$\%. At $T=20$, its zero-mean constraint has median residual $0.010$. Using the initial norm instead gives factors within $10$\% for $96$\% of checkpoints and horizons. The main exception is Orthogonal on cartpole, whose analytic radius exceeds the initial norm by a median $30$\% at $T=20$.

\paragraph{Ranking uncertainty.} Table~\ref{tab:pairs} includes unresolved comparisons. With ten seeds, bootstrap intervals collapse when one module wins every seed pair. Exact Mann--Whitney tests at the 5\% level resolve the same pairs for all six probes. The pendulum reversals at $T=5$ do not recur at $T=10$.

\input{tables/table_pairs.tex}

\paragraph{Verdicts.} Table~\ref{tab:pairs_main} lists the orderings on which the re-projected and the normalized probe agree, with every other probe and the protocol's verdict. Table~\ref{tab:verdicts} gives the verdict on every module pair in both designs, Tables~\ref{tab:gates} and~\ref{tab:gates_shared} the values behind step 3 and the support shares, and Table~\ref{tab:sensitivity} the verdicts under other thresholds. Step 3 compares medians over checkpoints. With or without a 10\% margin, it rejects $12$ of the eighteen base-model pairs, and no variant confirms a ranking in either design.

\input{tables/table_pairs_main.tex}

\input{tables/table_verdicts.tex}

\input{tables/table_gates.tex}

\input{tables/table_gates_shared.tex}

\input{tables/table_sensitivity.tex}

\paragraph{Models of known accuracy.}
\label{app:known}
On the encoder of each Non-Linear base model we build models whose accuracy is known. For each scoring episode, a model of rank $k$ predicts the encoding of a frame from another scoring episode whose state, standardized per coordinate, is among the $k$-th to $2k$-th nearest to the true state five steps ahead, drawn anew for each seed. The encoding is multiplied by the model's drift and a log-normal jitter. A smaller rank gives a smaller known state error. The model's null rollouts apply the same rule to the state reached from the nearest start of another episode (actions shuffled), to the next episode's state (starts swapped) and to the start state (untrained, near the identity), and the probe fitted on predictions sees the same model on the reference episodes. The most accurate model, of rank one at three times the reference radius, is paired with models of rank one, two, four, eight and $32$ at half that drift, each with ten seeds, and the protocol reads every pair as it reads the trained modules (Table~\ref{tab:known} gives the counts over all tasks and Table~\ref{tab:known_full} per task).

\input{tables/table_known_full.tex}


\paragraph{Training the probe on rescaled encodings.} Fitting the ordinary probe on encodings rescaled by random factors between a tenth and ten lowers median held-out $R^2$ from $0.85$ to $0.55$. At $T=5$, its error is $1.24$ times that of the re-projected ordinary probe. 

\paragraph{Removing encoder normalization.} Without final encoder normalization, adding norm to a direction probe improves held-out $R^2$ by at most $0.010$. The free modules' pooled re-projection factor remains low ($0.005$ at $T=20$). Re-projection can also raise the error. Orthogonal on cartpole shrinks to $0.21$ of its reference radius, and rescaling increases its error $9.8$-fold at $T=20$.

\subsection{Alternative probes and recurrent rescaling}
\label{app:alternative_readouts}
Pairwise comparisons refer to the seven orderings resolved by both the re-projected and normalized probes (Table~\ref{tab:pairs_main}).

\paragraph{The normalized probe.} The ordinary probe's architecture, fitted and evaluated on unit directions. It cannot respond to the norm, but it also cannot use any state the norm carries.

\paragraph{Fitting the probe on predictions.}
We fit a separate probe at each horizon on predictions from reference episodes. It confirms three of the seven orderings, reverses none, and leaves four unresolved. It shows how much of the future state the predictions contain, even when they lie away from the encodings.

On cartpole at $T=5$, probes fitted on predictions give similar relative errors for all architectures ($0.026$--$0.028$). Probes on the initial encoding and on untrained linear rollouts remain near the reference mean. The untrained Non-Linear network already reaches $0.045$, so this control distinguishes learning only for the linear modules. At $T=20$, prediction-fitted probes outperform the mean on $119$ of $120$ checkpoints, but start-encoding probes and untrained modules also often do so.

\paragraph{The neighbor probe.}
The neighbor probe uses the first probe fit's reference encodings and achieves median held-out $R^2=$ $0.86$. Its resolved pairwise orderings are unchanged with one, ten, or twenty neighbors. Disagreement with the re-projected probe persists on cartpole episodes supported by both modules under the held-out reference threshold. On these episodes, the re-projected probe favors both linear modules over Orthogonal, while neighbors and retrieval favor Orthogonal. All three favor Non-Linear over Orthogonal. Full-sample intervals are in Table~\ref{tab:pairs}.

\paragraph{Target retrieval.} Retrieval asks whether a prediction points closer to its own future encoding than to those of other episodes. Each prediction ranks its own target among the $100$ encoded scoring targets at the same horizon by cosine similarity, and rank $1$ is best. It needs no probe and so cannot extrapolate. It measures closeness in the encoder's geometry, which can differ from error in state units. Retrieval and the neighbor probe jointly resolve seven module pairs and agree on all of them.

\paragraph{Null rollouts.}
For action shuffling, we independently permute actions across scoring episodes at each step. Start swapping uses the preceding episode's initial encoding. Replaying checkpoints reproduces stored rollouts to relative error $1.9\times10^{-4}$. Table~\ref{tab:controls_full} reports all null-rollout results, including the untrained modules and the start-encoding probe.

\input{tables/table_controls.tex}

\input{tables/table_controls_full.tex}

\paragraph{Rescaling every rollout step.}
Re-projection changes only the final probe input. Norm drift can also alter subsequent dynamics: for $F_\theta(c\z,\mathbf a)=c\K\z+\B\mathbf a$, changing $c$ changes the latent's weight relative to the action term. We therefore compare final-input rescaling with rescaling every recurrent step, without retraining.

The median error ratio of recurrent to final-input rescaling is $1.00$ at $T=20$ for the trained modules and $0.97$--$1.07$ for TCRL, tested through $499$ steps. For Non-Linear, recurrent rescaling improves target alignment on only $15$ of $30$ checkpoints.
